%% ma: L11-B09-0006
%% lop: 11
%% chuong: 3
%% bai: 9
%% dang: 11.9.4
%% loai: TN4
%% muc: TH
%% dap_an: "B"
%% nguon: "chuyen-de-cac-so-dac-trung-cua-mau-so-lieu-khong-ghep-nhom-toan-10-chuong-trinh-moi.pdf (GV. Phan Nhật Linh) (P02), Bài 02 Dạng 1 Phần I Câu 36"
%% kien_thuc: "Lớp 11 Bài 9 (số trung bình của mẫu số liệu ghép nhóm)"
%% trang_thai: chinh-thuc
%% ghi_chu: "Mẫu số liệu ghép nhóm (vượt Lớp 10 Bài 13), đã chuyển sang Lớp 11 Bài 9. Tự duyệt theo yêu cầu giáo viên 22/10/2026."
\begin{ex}
Điều tra chiều cao (đơn vị: cm) của $100$ học sinh, ta có kết quả sau:
\begin{center}
\begin{tabular}{|c|c|c|c|c|c|c|}\hline
Chiều cao & $[150;152)$ & $[152;154)$ & $[154;156)$ & $[156;158)$ & $[158;160)$ & $[160;162)$\\\hline
Tần suất (\%) & $5$ & $18$ & $40$ & $26$ & $8$ & $3$\\\hline
\end{tabular}
\end{center}
Số trung bình của mẫu số liệu là
\choice{$155{,}36$}{\True $155{,}46$}{$154{,}98$}{$154{,}75$}
\loigiai{Giá trị đại diện: $151$; $153$; $155$; $157$; $159$; $161$.

$\overline x=151\cdot0{,}05+153\cdot0{,}18+155\cdot0{,}4+157\cdot0{,}26+159\cdot0{,}08+161\cdot0{,}03=155{,}46$ (cm). Chọn B.}
\end{ex}
